%%%%%%%%%%%%%%%%%%%%%%%%%%%%%%%%%%%%%%%%%%%%%%%%%%%%%%%%%%%%%%%%%%%%%%
% Lớp 8 | Chương X. Một số hình khối trong thực tiễn | Hình chóp tứ giác đều
% Dạng 8H0?39-2: Diện tích xung quanh, diện tích toàn phần, thể tích
%%%%%%%%%%%%%%%%%%%%%%%%%%%%%%%%%%%%%%%%%%%%%%%%%%%%%%%%%%%%%%%%%%%%%%


% nguồn: De_AI_Toan8_cuoi_hoc_ki_i_1 (nhập 21/11/2026)
\begin{bt}%[8H0H39-2]%[ma:6097d534]
	Cho hình chóp tứ giác đều $S.ABCD$ có cạnh đáy bằng $6\text{ cm}$ và chiều cao $SO = 4\text{ cm}$.
a) Tính độ dài trung đoạn của hình chóp.
b) Tính diện tích xung quanh, diện tích toàn phần và thể tích của hình chóp $S.ABCD$.
	\loigiai{
		a) Gọi $H$ là trung điểm của $CD$. Ta có $OH = \dfrac{AB}{2} = \dfrac{6}{2} = 3\text{ cm}$ và tam giác $SOH$ vuông tại $O$. \\
Áp dụng định lí Pythagore trong tam giác vuông $SOH$, độ dài trung đoạn $SH$ là: \\
$SH = \sqrt{SO^2 + OH^2} = \sqrt{4^2 + 3^2} = \sqrt{25} = 5\text{ cm}$. \\
b) Nửa chu vi đáy là $p = \dfrac{4 \cdot 6}{2} = 12\text{ cm}$. \\
Diện tích xung quanh: $S_{xq} = p \cdot SH = 12 \cdot 5 = 60\text{ cm}^2$. \\
Diện tích mặt đáy: $S_{\text{đáy}} = 6^2 = 36\text{ cm}^2$. Diện tích toàn phần: $S_{tp} = 60 + 36 = 96\text{ cm}^2$. \\
Thể tích hình chóp: $V = \dfrac{1}{3} S_{\text{đáy}} \cdot SO = \dfrac{1}{3} \cdot 36 \cdot 4 = 48\text{ cm}^3$.
	}
\end{bt}

% nguồn: De_AI_Toan8_cuoi_hoc_ki_i_1 (nhập 21/11/2026)
\begin{bt}%[8H0H39-2]%[ma:b4e6ba7c]
	Cho hình chóp tứ giác đều có cạnh đáy bằng $10\text{ cm}$ và chiều cao bằng $12\text{ cm}$.
a) Tính độ dài trung đoạn của hình chóp.
b) Tính diện tích xung quanh, diện tích toàn phần và thể tích của hình chóp.
	\loigiai{
		a) Khoảng cách từ tâm đáy đến một cạnh đáy là $r = \dfrac{10}{2} = 5\text{ cm}$. \\
Áp dụng định lí Pythagore, độ dài trung đoạn là $d = \sqrt{12^2 + 5^2} = \sqrt{144 + 25} = \sqrt{169} = 13\text{ cm}$. \\
b) Nửa chu vi đáy là $p = 2 \cdot 10 = 20\text{ cm}$. \\
Diện tích xung quanh: $S_{xq} = 20 \cdot 13 = 260\text{ cm}^2$. \\
Diện tích đáy: $S_{\text{đáy}} = 10^2 = 100\text{ cm}^2 \Rightarrow S_{tp} = 260 + 100 = 360\text{ cm}^2$. \\
Thể tích hình chóp: $V = \dfrac{1}{3} \cdot 100 \cdot 12 = 400\text{ cm}^3$.
	}
\end{bt}

% nguồn: De_AI_Toan8_cuoi_hoc_ki_i_1 (nhập 21/11/2026)
\begin{bt}%[8H0H39-2]%[ma:ef6c3aaa]
	Cho hình chóp tứ giác đều có cạnh đáy bằng $12\text{ cm}$ và độ dài trung đoạn bằng $10\text{ cm}$.
a) Tính chiều cao của hình chóp.
b) Tính diện tích xung quanh và thể tích của hình chóp.
	\loigiai{
		a) Khoảng cách từ tâm đáy đến trung điểm một cạnh đáy là $r = \dfrac{12}{2} = 6\text{ cm}$. \\
Áp dụng định lí Pythagore, chiều cao của hình chóp là $h = \sqrt{10^2 - 6^2} = \sqrt{100 - 36} = \sqrt{64} = 8\text{ cm}$. \\
b) Nửa chu vi đáy là $p = 2 \cdot 12 = 24\text{ cm}$, diện tích đáy là $S_{\text{đáy}} = 12^2 = 144\text{ cm}^2$. \\
Diện tích xung quanh: $S_{xq} = 24 \cdot 10 = 240\text{ cm}^2$. \\
Thể tích hình chóp: $V = \dfrac{1}{3} \cdot 144 \cdot 8 = 384\text{ cm}^3$.
	}
\end{bt}

% nguồn: De_AI_Toan8_cuoi_hoc_ki_i_1 (nhập 21/11/2026)
\begin{bt}%[8H0H39-2]%[ma:49130971]
	Cho hình chóp tứ giác đều có cạnh đáy bằng $16\text{ cm}$ và độ dài trung đoạn bằng $17\text{ cm}$.
a) Tính chiều cao của hình chóp.
b) Tính diện tích toàn phần và thể tích của hình chóp.
	\loigiai{
		a) Khoảng cách từ tâm hình vuông đáy đến cạnh đáy là $r = \dfrac{16}{2} = 8\text{ cm}$. \\
Áp dụng định lí Pythagore, chiều cao của hình chóp là $h = \sqrt{17^2 - 8^2} = \sqrt{289 - 64} = \sqrt{225} = 15\text{ cm}$. \\
b) Diện tích xung quanh: $S_{xq} = (2 \cdot 16) \cdot 17 = 32 \cdot 17 = 544\text{ cm}^2$. \\
Diện tích đáy: $S_{\text{đáy}} = 16^2 = 256\text{ cm}^2$. \\
Diện tích toàn phần: $S_{tp} = 544 + 256 = 800\text{ cm}^2$. \\
Thể tích hình chóp: $V = \dfrac{1}{3} \cdot 256 \cdot 15 = 1280\text{ cm}^3$.
	}
\end{bt}

% nguồn: De_AI_Toan8_cuoi_hoc_ki_i_1 (nhập 21/11/2026)
\begin{bt}%[8H0H39-2]%[ma:8207ad4c]
	Cho hình chóp tứ giác đều $S.ABCD$ có cạnh đáy $AB = 12\text{ cm}$ và cạnh bên $SA = 10\text{ cm}$.
a) Tính độ dài trung đoạn của hình chóp.
b) Tính diện tích xung quanh và diện tích toàn phần của hình chóp.
	\loigiai{
		a) Gọi $H$ là trung điểm của $AB$. Vì tam giác $SAB$ cân tại $S$ nên $SH \perp AB$ và $AH = \dfrac{AB}{2} = 6\text{ cm}$. \\
Áp dụng định lí Pythagore trong tam giác $SAH$ vuông tại $H$, độ dài trung đoạn là: \\
$SH = \sqrt{SA^2 - AH^2} = \sqrt{10^2 - 6^2} = \sqrt{64} = 8\text{ cm}$. \\
b) Diện tích xung quanh: $S_{xq} = (2 \cdot 12) \cdot 8 = 192\text{ cm}^2$. \\
Diện tích đáy: $S_{\text{đáy}} = 12^2 = 144\text{ cm}^2$. \\
Diện tích toàn phần: $S_{tp} = 192 + 144 = 336\text{ cm}^2$.
	}
\end{bt}

% nguồn: De_AI_Toan8_cuoi_hoc_ki_i_1 (nhập 21/11/2026)
\begin{bt}%[8H0H39-2]%[ma:3f389bcf]
	Cho hình chóp tứ giác đều $S.ABCD$ có cạnh đáy bằng $10\text{ cm}$ và cạnh bên bằng $13\text{ cm}$.
a) Tính độ dài trung đoạn của hình chóp.
b) Tính diện tích xung quanh và diện tích toàn phần của hình chóp.
	\loigiai{
		a) Gọi $M$ là trung điểm cạnh $BC$. Trong tam giác cân $SBC$, trung đoạn $SM \perp BC$ và $BM = \dfrac{10}{2} = 5\text{ cm}$. \\
Áp dụng định lí Pythagore trong tam giác $SBM$ vuông tại $M$: \\
$SM = \sqrt{SB^2 - BM^2} = \sqrt{13^2 - 5^2} = \sqrt{144} = 12\text{ cm}$. \\
b) Diện tích xung quanh: $S_{xq} = (2 \cdot 10) \cdot 12 = 240\text{ cm}^2$. \\
Diện tích toàn phần: $S_{tp} = S_{xq} + 10^2 = 240 + 100 = 340\text{ cm}^2$.
	}
\end{bt}

% nguồn: De_AI_Toan8_cuoi_hoc_ki_i_1 (nhập 21/11/2026)
\begin{bt}%[8H0H39-2]%[ma:604797f7]
	Một hình chóp tứ giác đều có diện tích xung quanh bằng $180\text{ cm}^2$ và độ dài trung đoạn bằng $9\text{ cm}$.
a) Tính độ dài cạnh đáy của hình chóp.
b) Tính diện tích toàn phần của hình chóp.
	\loigiai{
		a) Từ công thức $S_{xq} = p \cdot d$, ta tính được nửa chu vi đáy là $p = \dfrac{S_{xq}}{d} = \dfrac{180}{9} = 20\text{ cm}$. \\
Chu vi đáy hình vuông là $C = 2p = 40\text{ cm}$, suy ra độ dài cạnh đáy là $a = \dfrac{40}{4} = 10\text{ cm}$. \\
b) Diện tích đáy là $S_{\text{đáy}} = 10^2 = 100\text{ cm}^2$. \\
Diện tích toàn phần của hình chóp là $S_{tp} = 180 + 100 = 280\text{ cm}^2$.
	}
\end{bt}

% nguồn: De_AI_Toan8_cuoi_hoc_ki_i_1 (nhập 21/11/2026)
\begin{bt}%[8H0H39-2]%[ma:22c36f66]
	Một hình chóp tứ giác đều có thể tích bằng $192\text{ cm}^3$ và chiều cao bằng $9\text{ cm}$.
a) Tính diện tích mặt đáy của hình chóp.
b) Tính độ dài cạnh đáy và chu vi mặt đáy của hình chóp.
	\loigiai{
		a) Từ công thức $V = \dfrac{1}{3} S_{\text{đáy}} \cdot h$, ta có: \\
$S_{\text{đáy}} = \dfrac{3V}{h} = \dfrac{3 \cdot 192}{9} = 64\text{ cm}^2$. \\
b) Vì mặt đáy là hình vuông có diện tích $64\text{ cm}^2$ nên độ dài cạnh đáy là $a = \sqrt{64} = 8\text{ cm}$. \\
Chu vi mặt đáy là $4 \cdot 8 = 32\text{ cm}$.
	}
\end{bt}

% nguồn: De_AI_Toan8_cuoi_hoc_ki_i_1 (nhập 21/11/2026)
\begin{bt}%[8H0H39-2]%[ma:e2d86e1f]
	Một hình chóp tứ giác đều có thể tích bằng $400\text{ cm}^3$ và cạnh đáy bằng $10\text{ cm}$.
a) Tính chiều cao của hình chóp.
b) Tính độ dài trung đoạn và diện tích xung quanh của hình chóp.
	\loigiai{
		a) Diện tích mặt đáy là $S_{\text{đáy}} = 10^2 = 100\text{ cm}^2$. \\
Chiều cao của hình chóp là $h = \dfrac{3V}{S_{\text{đáy}}} = \dfrac{3 \cdot 400}{100} = 12\text{ cm}$. \\
b) Khoảng cách từ tâm đáy đến cạnh đáy là $r = \dfrac{10}{2} = 5\text{ cm}$. \\
Độ dài trung đoạn là $d = \sqrt{12^2 + 5^2} = 13\text{ cm}$. \\
Diện tích xung quanh là $S_{xq} = (2 \cdot 10) \cdot 13 = 260\text{ cm}^2$.
	}
\end{bt}

% nguồn: De_AI_Toan8_cuoi_hoc_ki_i_1 (nhập 21/11/2026)
\begin{bt}%[8H0H39-2]%[ma:45002ed1]
	Một hình chóp tứ giác đều có chu vi mặt đáy bằng $24\text{ cm}$ và diện tích xung quanh bằng $60\text{ cm}^2$.
a) Tính cạnh đáy và độ dài trung đoạn của hình chóp.
b) Tính chiều cao và thể tích của hình chóp.
	\loigiai{
		a) Cạnh đáy hình vuông là $a = \dfrac{24}{4} = 6\text{ cm}$, nửa chu vi đáy là $p = \dfrac{24}{2} = 12\text{ cm}$. \\
Độ dài trung đoạn là $d = \dfrac{S_{xq}}{p} = \dfrac{60}{12} = 5\text{ cm}$. \\
b) Khoảng cách từ tâm đáy đến cạnh đáy là $r = \dfrac{6}{2} = 3\text{ cm}$. \\
Chiều cao hình chóp là $h = \sqrt{5^2 - 3^2} = 4\text{ cm}$. \\
Thể tích hình chóp là $V = \dfrac{1}{3} \cdot 6^2 \cdot 4 = 48\text{ cm}^3$.
	}
\end{bt}

% nguồn: De_AI_Toan8_cuoi_hoc_ki_i_1 (nhập 21/11/2026)
\begin{bt}%[8H0H39-2]%[ma:9b495d6a]
	Cho hình chóp tam giác đều $S.ABC$ có cạnh đáy $AB = 10\text{ cm}$ và cạnh bên $SA = 13\text{ cm}$.
a) Tính độ dài trung đoạn của hình chóp.
b) Tính diện tích xung quanh của hình chóp $S.ABC$.
	\loigiai{
		a) Gọi $H$ là trung điểm cạnh $BC$. Vì tam giác $SBC$ cân tại $S$ nên $SH \perp BC$ và $BH = \dfrac{BC}{2} = 5\text{ cm}$. \\
Áp dụng định lí Pythagore trong tam giác $SBH$ vuông tại $H$, độ dài trung đoạn là: \\
$SH = \sqrt{SB^2 - BH^2} = \sqrt{13^2 - 5^2} = \sqrt{144} = 12\text{ cm}$. \\
b) Nửa chu vi đáy là $p = \dfrac{3 \cdot 10}{2} = 15\text{ cm}$. \\
Diện tích xung quanh của hình chóp là $S_{xq} = 15 \cdot 12 = 180\text{ cm}^2$.
	}
\end{bt}

% nguồn: De_AI_Toan8_cuoi_hoc_ki_i_1 (nhập 21/11/2026)
\begin{bt}%[8H0H39-2]%[ma:1f67ea8f]
	Cho hình chóp tam giác đều $S.ABC$ có cạnh đáy bằng $12\text{ cm}$ và cạnh bên bằng $10\text{ cm}$.
a) Tính độ dài trung đoạn của hình chóp.
b) Tính diện tích một mặt bên và diện tích xung quanh của hình chóp.
	\loigiai{
		a) Gọi $I$ là trung điểm của cạnh $BC$. Ta có $BI = \dfrac{12}{2} = 6\text{ cm}$ và $SI \perp BC$. \\
Áp dụng định lí Pythagore trong tam giác $SBI$ vuông tại $I$: $SI = \sqrt{10^2 - 6^2} = \sqrt{64} = 8\text{ cm}$. \\
b) Diện tích một mặt bên là $S_{SBC} = \dfrac{1}{2} \cdot BC \cdot SI = \dfrac{1}{2} \cdot 12 \cdot 8 = 48\text{ cm}^2$. \\
Diện tích xung quanh của hình chóp là $S_{xq} = 3 \cdot 48 = 144\text{ cm}^2$.
	}
\end{bt}

% nguồn: chuyen_dang (nhập 22/11/2026)
\begin{bt}%[8H0V39-2]%[ma:15b20df6]
% nguon: Đề tham khảo Tuyển sinh 10 TP.HCM 2025-2026 – Quận 3 – Đề 2 – Bài 4
Một hình chóp tứ giác đều có chiều dài cạnh đáy là $x\text{ (cm)}$ và chiều dài đường cao của mặt bên (trung đoạn) ngắn hơn cạnh đáy $2\text{ cm}$.
\begin{enumerate}[a)]
	\item Viết biểu thức $S_{xq}$ biểu diễn diện tích xung quanh của hình chóp theo $x$.
	\item Biết diện tích toàn phần của hình chóp tứ giác đều này là $207\text{ cm}^2$, hãy tính độ dài cạnh đáy $x$.
\end{enumerate}
\loigiai{
\begin{enumerate}[a)]
	\item Chiều dài đường cao của mặt bên (trung đoạn của hình chóp) là: $d = x - 2\text{ (cm)}$ (với $x > 2$).\\
	Hình chóp tứ giác đều có $4$ mặt bên là các tam giác cân có đáy $x$ và đường cao $d = x - 2$.\\
	Biểu thức biểu diễn diện tích xung quanh của hình chóp là:
	$S_{xq} = 4 \cdot \left(\dfrac{1}{2} x \cdot (x - 2)\right) = 2x(x - 2) = 2x^2 - 4x\text{ (cm}^2\text{)}$.
	\item Diện tích đáy của hình chóp tứ giác đều là: $S_{\text{đáy}} = x^2\text{ (cm}^2\text{)}$.\\
	Diện tích toàn phần của hình chóp là:
	$S_{tp} = S_{xq} + S_{\text{đáy}} = (2x^2 - 4x) + x^2 = 3x^2 - 4x\text{ (cm}^2\text{)}$.\\
	Theo đề bài, $S_{tp} = 207\text{ cm}^2$, ta có phương trình:
	$3x^2 - 4x = 207 \Leftrightarrow 3x^2 - 4x - 207 = 0$.\\
	Ta có $\Delta' = (-2)^2 - 3\cdot(-207) = 4 + 621 = 625 = 25^2 > 0$.\\
	Vì $x > 2$ nên phương trình có nghiệm hợp lệ duy nhất:
	$x = \dfrac{2 + 25}{3} = \dfrac{27}{3} = 9\text{ (cm)}$.\\
	Vậy độ dài cạnh đáy của hình chóp tứ giác đều là $9\text{ cm}$.
\end{enumerate}
}
\end{bt}
